Ion thrusters have emerged as a primary propulsion choice for long-duration space missions. Their performance is heavily dependent on the Power Processing Unit (PPU) and its ability to efficiently step up low bus voltages to kilovolt levels. This paper evaluates the trade-offs between hard-switched and soft-switched (resonant) DC-DC converter topologies in terms of efficiency, mass, and reliability in the context of emerging power electronic semiconductors, such as gallium arsenide (GaAs) and silicon carbide (SiC).

While hard-switched architectures inherently suffer from their switching losses being proportional to the operating frequency, and thus limited power density, resonant topologies utilize Zero Voltage Switching (ZVS) and Zero Current Switching (ZCS) to decouple efficiency from frequency. Our review finds that while soft-switching can achieve efficiencies up to 99~\% and reduce electromagnetic interference, it introduces additional mass by advanced circuit complexity. The analysis concludes that modern PPU design must balance the superior high-frequency efficiency of resonant converters against the increased component count and potential reliability vulnerabilities inherent in the complex architectures.